\pagestyle{fancy}\thispagestyle{plain}

\rule{6.8in}{1pt}\vskip.1in

{\large Assignment 12: \csname topic12\endcsname}\smallskip

\rule{6.8in}{1pt}\vskip.1in

\textit{Due by  \HWdueTime,  eastern, on \dueDay, \csname dateWeek12\endcsname}

\rule{6.8in}{1pt}\vskip.1in

{\bf\textit{Suggested readings  for this problem set:}}
\begin{itemize}[itemsep=1pt]
\item Section 4.2, p.~139-144 (stop at the proof of Theorem 4.2.6)
\end{itemize}

All readings are from \bookAuthors, \emph{\bookName}.

\rule{6.8in}{1pt}\vskip.1in

{\bf\textit{Extra Practice Problems}} from \emph{\bookName} (not to turn in):

\begin{itemize}[itemsep=1pt]
\item With answers or hints (in the back of the book): Section 4.2 \#1(a), 3(ad), 4(a), 5(a), 12(a)
\item Without answers: Section 4.2 \#1, 3, 4  
\item \href{https://dmzb.github.io/teaching/2024Spring220/handouts/220-H14-equivalence-relations.pdf}{Handout 14}

\end{itemize}

\rule{6.8in}{1pt}\vskip.1in

 {\bf\textit{Assignment:}} due \dueDay, \csname dateWeek12\endcsname, \HWdueTime, via Gradescope (\gradescopeCode):

\rule{6.8in}{1pt}\vskip.1in

\begin{enumerate}
\item Define a relation $R$ on the set $\mathbb{R}$ as follows: we say that $x \sim y$ if $|x| = |y|$. (Recall that $|x|$ is the absolute value of $x$) Determine whether this relation is reflexive, symmetric, transitive (and in each case give a proof or disproof).

\item Let $A = \{1,2,3\}$ and define a relation on $A$ by $a \sim b$ if $a + b \neq 3$. Determine whether this relation is reflexive, symmetric, transitive (and in each case give a proof or disproof).

\item Define a relation on $\mathbf{Z}$ given by $a \sim b$ if $a-b$ is divisible by $3$.
 \begin{enumerate}
 \item Prove that this is an equivalence relation.
 \item What integers are in the equivalence class of 18? (No proof necessary.)
 \item What integers are in the equivalence class of 31? (No proof necessary.) 
 \item How many distinct equivalence classes are there? What are they? (No proof necessary.)
 \end{enumerate}
\item Define a relation on $\mathbf{Z}$ given by $a \sim b$ if $a^2-b^2$ is divisible by $4$.
 \begin{enumerate}
 \item Prove that this is an equivalence relation.
 \item How many distinct equivalence classes are there? What are they? (No proof necessary.)
 \end{enumerate} 
\item Let $A$ be a set, and let $P(A)$ be the power set of $A$. Assume that $A$ is not the empty set. Define a relation on $P(A)$ by $X \sim Y$ if $X \subseteq Y$. Is this relation reflexive, symmetric, and/or transitive? In each case give a proof, or disprove with a counterexample. (For a counterexample, give an example of $A$, $X$, and $Y$ that disproves the statement.)

\item Define a relation $R$ on the set $\mathbb{R}$ as follows: we say that $x \sim y$ if $x - y \in \mathbb{Q}$. Determine whether this relation is reflexive, symmetric, transitive (and in each case give a proof or disproof).
  
\item Define a relation $R$ on the set $\mathbb{Z}$ as follows: we say that $x \sim y$ if there exists a non-negative integer $n$ such that $x = 2^ny$. Determine whether this relation is reflexive, symmetric, transitive (and in each case give a proof or disproof).

\item Define a relation $R$ on the set $\mathbb{Z}\times \mathbb{Z} - \{(0,0)\}$ as follows: we say that $(a,b) \sim (c,d)$ if $ad = bc$. Prove that this is an equivalence relation. What are the equivalence classes?
 \begin{enumerate}
 \item Prove that this is an equivalence relation.
 \item What pairs $(c,d)$ are in the equivalence class of $(1,1)$? (No proof necessary.)
 \item What pairs $(c,d)$ are in the equivalence class of $(1,2)$? (No proof necessary.)   
 \item How many distinct equivalence classes are there? What are they? (No proof necessary.)
 \end{enumerate} 
  
\end{enumerate}




%%% Local Variables: 
%%% mode: latex
%%% TeX-master: "../assignments-S24-220"
%%% End: 